\section{Semiconductor Photodetector Theory}

\subsection{Semiconductor Current Density}
\label{sec:current_density}

Let $j_{\mathrm{n}}$ be the electron current density and $j_{\mathrm{p}}$ be the
hole current density. The total current density can be calculated as
$j_{\mathrm{tot}}= j_{\mathrm{n}}(z) + j_{\mathrm{p}}(z)$. In a device with
constant cross section $A$, the total current is obtained as $I=A j_{\mathrm{tot}}$.

In semiconductor physics, currents originate either from an electric field (drift
current) or from a concentration gradient (diffusion current). For nondegenerate
semiconductors, this results in
\begin{align}
	j_{\mathrm{p}} & = e \mu_{\mathrm{p}} p E - e D_{\mathrm{p}} \partial_{z}p,
	\label{eq:j_p}                                                             \\
	j_{\mathrm{n}} & = -e \mu_{\mathrm{n}} n E - e D_{\mathrm{n}} \partial_{z}n.
	\label{eq:j_n}
\end{align}

Those equations feature
\begin{itemize}
	\item Elementary charge $e$
	\item Charge carrier concentration $n / p$
	\item Mobility of the respective charge carrier $\mu_{\mathrm{n/p}}$
	\item Electric field $E$
	\item The Einstein coefficient $D= kT \mu/q$, with the Boltzmann constant
		$k$, the temperature $T$ and the charge $q$ of the respective carrier
\end{itemize}

\subsection{Electron-Hole Continuity Equation}
\label{sec:continuity}

In an infinitesimally small volume at $\mathbf{r}$, the temporal change of electron
and hole density is determined by generation rates, summarized in
$G_{\mathrm{n/p}}(\mathbf{r})$ for electrons / holes, recombination rates,
summarized in $U_{\mathrm{n/p}}$ for electrons / holes and the current densities
$\mathbf{j}_{\mathrm{n/p}}$ that flow through the volume's surface. The elementary
charge is denoted by $e$.
\begin{align}
	\partial_{t}n & = G_{\mathrm{n}}- U_{\mathrm{n}}
	+ \frac{1}{e}\nabla \cdot \mathbf{j}_{\mathrm{n}},                     \\
	\partial_{t}p & = G_{\mathrm{p}}- U_{\mathrm{p}}
	- \frac{1}{e}\nabla \cdot \mathbf{j}_{\mathrm{p}}.
\end{align}
In most two-terminal devices, characteristics only depend on one direction $(z)$,
while being independent in its cross section. This reduces the complexity as
\begin{align}
	\partial_{t}n(z) & = G_{\mathrm{n}}(z) - U_{\mathrm{n}}(z)
	+ \frac{1}{e}\partial_{z}j_{\mathrm{n}}(z),                \\
	\partial_{t}p(z) & = G_{\mathrm{p}}(z) - U_{\mathrm{p}}(z)
	- \frac{1}{e}\partial_{z}j_{\mathrm{p}}(z).
	\label{eq:continuity_1d}
\end{align}

\subsection{Photon Continuity Equation}
\label{sec:photon_continuity}

Let $\Phi(z)$ denote the photon flux, $n_{\mathrm{ph}}(z)$ the photon density and
$U_{\mathrm{ph}}$ the photon absorption rate per unit volume in a two-terminal
device. For a small slab $[z, z+\delta]$ of cross section $A$, the following
statement is true:
\begin{equation}
	\underbrace{ \partial_{t}(n_{\mathrm{ph}}(z,t)A\delta) }_{ \text{change of photons} }
	= \underbrace{ \Phi(z)A }_{ \substack{\text{incoming} \\ \text{photons}} }
	- \underbrace{ \Phi(z+\delta)A }_{ \substack{\text{outgoing} \\ \text{photons}} }
	- \underbrace{ U_{\mathrm{ph}}A\delta }_{ \substack{\text{photon} \\ \text{sink}} }.
\end{equation}
Dividing by $\delta A$ and $\delta \to 0$ yields the photon continuity equation

\begin{equation}
	\partial_{t}n_{\mathrm{ph}}
	= - \frac{\mathrm{d} \Phi}{\mathrm{d}z}-U_{\mathrm{ph}}.
\end{equation}
In steady state $\partial_{t}n_{\mathrm{ph}}= 0$, resulting in
\begin{equation}
	U_{\mathrm{ph}}=-\frac{\mathrm{d} \Phi}{\mathrm{d}z}.
	\label{eq:photon_continuity_steady}
\end{equation}

\subsection{Lambert-Beer Law}
\label{sec:lambert_beer}

Let a beam of light enter a material and let the $z$ direction be parallel to the
beam direction. Divide the material into thin slices of width $\delta$. The radiant
flux that is absorbed in the slice is proportional to the incoming flux and the
slice width with proportionality constant $\alpha$:
\begin{equation}
	\Phi(z) - \Phi(z+\delta) = \alpha(h\nu, z) \Phi(z) \delta.
\end{equation}
Approaching $\delta \to 0$ yields the following differential equation:
\begin{equation}
	\frac{\mathrm{d}\Phi(z)}{\mathrm{d}z}=-\alpha(h\nu, z) \Phi(z).
\end{equation}
The quantity $\alpha(h\nu, z)$ is the absorption coefficient, which depends on the
photon energy $h\nu$ and on the position $z$. A general solution for
$z$-dependent $\alpha$ can be constructed using a separation of variables ansatz:
\begin{align}
	\frac{\Phi}{\Phi_{0}}              & = \exp\left( -\int_{0}^{z}\alpha(h \nu, z')
	\, \mathrm{d}z' \right),                                                         \\
	\frac{\mathrm{d}\Phi}{\mathrm{d}z} & = -\alpha(h\nu, z) \Phi_{0}
	\exp\left( -\int_{0}^{z}\alpha(h\nu, z') \, \mathrm{d}z' \right).
	\label{eq:flux_general}
\end{align}
Here, $\Phi_{0}=\Phi(z=0)$ is the incident flux entering the material. If $\alpha$
is independent of $z$, this reduces to the well-known Lambert-Beer law
\begin{equation}
	\Phi(z) = \Phi_{0}\exp(-\alpha(h\nu) z).
	\label{eq:lambert_beer}
\end{equation}
